% author: Matheus T. dos Santos
% co-authors:
%   - Claude (claude-opus-5), 2026-08-03 -- revisao mecanica autorizada por Matheus nesta data:
%     ortografia/acentuacao, concordancia, pontuacao, aspas LaTeX (``''), labels/refs invalidos
%     e unificacao \cite -> \citep. Nenhum conteudo, argumento ou frase autoral foi reescrito.
%   - Claude (claude-opus-4-8), 2026-08-07 -- padronizacao de genero de HAL (feminino)
%     autorizada por Matheus nesta data (suspensao pontual da Regra 3). Nenhum conteudo
%     ou frase autoral foi alterado.
%   - Claude (claude-opus-5), 2026-08-10 -- troca de chave de citacao autorizada por Matheus
%     nesta data (suspensao pontual da Regra 3): \citep{book:deadline-requirement} ->
%     \citep{article:liu-layland-hard-real-time} na secao 1.1, por falta de acesso ao livro do
%     Buttazzo. Somente a chave mudou; nenhuma palavra do texto foi tocada.
\mychapter{Introdução}{cap:introducao}

\section{Contextualização}

Máquinas que necessitam de controle automático são uma grande inovação da modernidade. Essas máquinas estão presentes não só em fábricas e indústrias, como motores e bombas, como também nos produtos mais sutis, como um marca-passo para sincronizar as batidas do coração. Controlar essas máquinas indica executar algum processo para que elas sigam da melhor forma um valor de referência. Isso pode ser feito de maneira analógica ou digital. Com a diminuição de custo e popularização da computação, utilizar a teoria de controle \citep{book:ogata} em computadores, na forma de programas (ou algoritmos), se tornou cada vez mais viável e desejado. Além disso, o desenvolvimento e evolução de processadores \ac{RISC}, como ARM e RISC-V, permitiram um avanço ainda maior na implantação destes algoritmos de controle em dispositivos mais baratos e com periféricos para interagir com o mundo exterior. Esses dispositivos são os \ac{SoC}: processadores com periféricos para interagir com o meio externo, como: \ac{GPIO}, \ac{UART}, \ac{SPI}, etc. Um tipo de \ac{SoC} é uma \ac{MCU}, na qual está contido o core (processador), a memória e os periféricos. É importante os microcontroladores possuírem periféricos de interação com o mundo exterior, pois os algoritmos de controle executam em malha fechada. Isso quer dizer que o algoritmo envia sinais e recebe sinais como resposta ao estado da planta que ele está controlando. Esse controle deve ser feito em um período fixo e o tempo máximo para atuar em cada ciclo é chamado de deadline. Perder deadline aqui é diferente de perder deadline em um Desktop e a interface gráfica travar. Perder deadline aqui é o risco de o sistema controlado se desestabilizar e causar prejuízos. Não só financeiros, como uma planta de uma fábrica parada, mas também pessoais, podendo um motor instável causar acidentes ou um marca-passo causar uma arritmia \citep{article:liu-layland-hard-real-time}.

Note que, no algoritmo de controle, o usuário pode definir a referência e isso chega para o algoritmo de controle de forma assíncrona. A tarefa de controle, por ser crítica e não poder perder deadline, é executada em uma tarefa dedicada, enquanto a interação com o usuário (uma \ac{IHM}, por exemplo), por ser de baixa prioridade, está em outra tarefa. Além disso, ainda existe a aquisição de dados (a resposta do estado do mundo exterior), obtida via interrupção. Ou seja, sempre que um novo dado está disponível, a \ac{MCU} é interrompida pelo periférico e este salva o valor lido dos sensores. Note que existem dois contextos fora da tarefa principal que contém o algoritmo de controle. E esses dois contextos precisam comunicar suas atualizações para a tarefa principal. Desta forma, existem estados que mais de um contexto precisa utilizar (ou para atualizar, ou para ser lido e utilizado) e acessar esses dados concorrentemente, sem sincronização, gera um data race. Perceba que é da natureza da aplicação possuir esses estados compartilhados e não um erro do programador.

O microcontrolador (ou \ac{SoC}) é barato por um motivo bem direto: ele possui recursos bem limitados. Apenas o necessário para fazer uma tarefa específica, que neste caso é controlar uma planta. Algumas centenas de kilobytes de memória e algumas centenas de MHz de processamento já são o suficiente para essa tarefa. Utilizar os gigabytes de memória e os GHz de configurações desktops seria gastar dinheiro a mais, pois uma configuração com menos recurso já resolve o problema. Para esse sistema com poucos recursos é fundamental utilizar uma linguagem que, além de permitir a gerência desses recursos, também permita que estes não sejam desperdiçados. Com isso, uma das linguagens que se encaixam neste requisito é a linguagem C. É nesta linguagem que as \acp{HAL} dos microcontroladores são construídas. Por ser uma linguagem utilizada desde a década de 70, ela possui uma abundância de código legado. Possuir o poder de gerenciar os recursos (como memória) é uma grande responsabilidade, e se não utilizado do jeito correto pode gerar problemas, tal como não sincronizar o acesso concorrente a um dado e gerar data races. Com isso, um padrão foi criado para diminuir os erros cometidos pelo desenvolvedor, o padrão \ac{MISRA} \citep{doc:MISRA}. Se o desenvolvedor seguir esse padrão, o seu código vai estar menos propenso aos erros citados. Porém, a própria frase denuncia o problema: ``\textbf{Se} o desenvolvedor seguir...''. Se o desenvolvedor não seguir, o problema continua lá. O \ac{MISRA} é apenas uma convenção, não impede a compilação de um programa com problemas. Além disso, a análise estática que verifica se o padrão \ac{MISRA} está sendo seguido não é completa. Para contornar isso, é utilizado um sanitizer \citep{doc:TSan}. Com o sanitizer é possível executar o programa e capturar problemas, como data races. Entretanto, os sanitizers apenas minimizam o problema, pois surgem outros com ele, pois: a cobertura é dinâmica, visto que só o que é exercitado durante a execução do programa é pego pelo sanitizer. Se um data race existir e o programa não executar aquele trecho, o sanitizer não o pega. Outro ponto é sua execução, pois: alguns sanitizers, como o \ac{TSan}, não executam em ambiente embarcado, apenas no desktop; e cada sanitizer identifica um tipo de bug, sendo necessário executar uma vez para cada tipo de sanitizer. O sanitizer apenas detecta o bug que foi exercitado, ele não previne qualquer bug daquela classe. Deste modo, a garantia hoje para sistemas de controle depende da disciplina do desenvolvedor e de o bug ter sido exercitado. Bugs latentes no código não são detectados e vão para a produção causando danos à planta controlada.

Uma linguagem mais recente, chamada Rust, desloca a garantia de detecção em tempo de execução para a recusa em tempo de compilação, através dos mecanismos: Send/Sync e borrow checker \citep{article:rust-safe-soundness} e \citep{article:rust-critical}. Além de mover a garantia, essa linguagem ainda mantém as funcionalidades de: controle sobre recursos (como memória) e baixa utilização de recursos. Apesar de ser uma linguagem recente, quando comparada com a linguagem C, ela já possui mais de 10 anos. A proposta de utilizá-la em algoritmos embarcados de controle somente agora se deve ao ecossistema de sistemas embarcados, da linguagem Rust, estar maduro recentemente, por exemplo: a maturidade do no\_std (bare-metal) da ESP32 ter ocorrido somente por volta de 2025 \citep{url:esp-rust-no-std-stable}. A linguagem não oferece uma bala de prata, visto que existem custos e incógnitas associados a mover a garantia: não saber onde os resquícios de unsafe (código fora das garantias da linguagem) se encontram nos padrões de data races e forçar sincronização. Essa sincronização tem um custo e esse custo pode afetar, ou não, deadlines em malhas rápidas, causando danos como em motores instáveis e em marca-passos fora do ritmo. Para investigar o custo das garantias de sincronização é necessário um veículo que permita a implementação de data races em algoritmos de controle realistas. Entretanto, o ecossistema Rust de controle não possui tal veículo que reúne suporte a: sistemas embarcados (no\_std), tratamento de concorrência e genérico pleno (\ref{sec:control-repos}). Com isso, a biblioteca Aule \citep{aule} foi construída, para ser veículo desta investigação.

Com essas lacunas apontadas, a pergunta que surge sobre utilizar Rust em algoritmos embarcados de controle é: onde se encontra a fronteira entre safe e unsafe do Rust em data races, no domínio de algoritmos embarcados de controle? Respondida a pergunta anterior, qual o custo cobrado pelas garantias da linguagem sobre data races?

\section{Objetivos}

\subsection{Objetivo Geral}

O objetivo geral é caracterizar a fronteira entre safe e unsafe dos data races presentes em algoritmos de controle que compartilham estado entre múltiplos contextos e avaliar o custo para implementar a sincronização imposta pelo lado safe da fronteira.

\subsection{Objetivos Específicos}

Os objetivos 1, 2 e 3 estão desenvolvidos na qualificação, enquanto os objetivos 4, 5, 6, 7 e 8 serão desenvolvidos na pós-qualificação.

\begin{enumerate}
\item Propor uma taxonomia de padrões de data races presentes em algoritmos de controle que compartilham estado entre múltiplos contextos;
\item Caracterizar a fronteira entre safe e unsafe destes padrões de data races;
\item Catalogar o espaço de design de implementar as garantias impostas para tornar os data races, presentes no lado safe da fronteira, inexprimíveis;
\end{enumerate}

\begin{enumerate}
\setcounter{enumi}{3}
\item Implementar, na biblioteca Aule, um subconjunto das garantias (um por eixo do espaço de design) impostas para tornar os data races, presentes no lado safe da fronteira, inexprimíveis;
\item Avaliar o overhead de tempo de execução destas implementações (em Rust) e em implementações em C, em controladores com realimentação de estados, em controle de um pêndulo invertido, em placas com ESP32 (Xtensa);
\item Avaliar a perda de deadlines nos algoritmos de controle implementados em Rust e em C;
\item Comparar o overhead de tempo de execução e a perda de deadlines entre a implementação em Rust e em C de controladores que compartilham estado entre múltiplos contextos;
\item Comparar a verificação em tipos do Rust, que torna os data races inexprimíveis, e o estado da arte em C+\allowbreak MISRA+\allowbreak sanitizers, que detecta os data races em tempo de execução.
\end{enumerate}

\section{Visão Geral da Qualificação}

Dado que a proposta está descrita, o restante é organizado assim: no capítulo \ref{cap:fundamentacao} está presente a fundamentação com as explicações e definições de data race e garantias do Rust. É neste capítulo que existe o conhecimento necessário para a compreensão do resto do texto. Na metodologia, capítulo \ref{cap:metodologia}, está presente a taxonomia, a fronteira entre safe e unsafe e o catálogo do espaço de design. No capítulo a seguir (\ref{cap:resultados}), são descritos os resultados parciais: o estado atual da biblioteca Aule; o caso demonstrativo P1; e o panorama do ecossistema open-source do Rust, para bibliotecas de controle. Por fim, no capítulo \ref{cap:cronograma} está descrito o cronograma pós-qualificação.

Na qualificação, os objetivos de 1 a 3 estão concluídos, enquanto os objetivos de 4 a 8 serão concluídos na dissertação. Nos resultados parciais está a evidência já colhida. A ordem dos capítulos é construída para facilitar o entendimento, pois se começa com a pergunta, levando para os objetivos do trabalho. Após isso são apresentados os fundamentos necessários para o entendimento geral, seguido do método utilizado para coletar e verificar as evidências, fechando com o cronograma do que será realizado após este trabalho.
